% !TEX program = latexmk
% !TEX options = -pdf -interaction=nonstopmode -halt-on-error
% ---------------------------------------------------------------------------
%  main.tex — Informe final de trabajo de grado (UPB).
%  Plantilla/esqueleto. Compile with:  ./build.sh
%
%  Language strategy:
%    • Front matter (portadas, contenido, resúmenes) is SPANISH and institutional.
%    • Body chapters are written in ENGLISH and translated to Spanish later.
%    • Institutional Spanish strings live ONLY in preamble/upb-strings.tex.
% ---------------------------------------------------------------------------

\documentclass[11pt, letterpaper]{article}

% ---------------------------------------------------------------------------
%  packages.tex — Shared packages for the UPB report template.
%  Load order matters: babel -> csquotes -> biblatex -> hyperref (hyperref last).
% ---------------------------------------------------------------------------

% --- Encoding & fonts -----------------------------------------------------
\usepackage[T1]{fontenc}
\usepackage[utf8]{inputenc}
\usepackage{tgtermes}             % TeX Gyre Termes — Times-like serif

% --- Page geometry: US Letter (612x792pt), ~2.5 cm margins ----------------
\usepackage[letterpaper, margin=2.5cm, headheight=14pt]{geometry}

% --- Languages -------------------------------------------------------------
% Spanish = front matter / institutional strings; English = body chapters.
% "main=english" makes English the default language of the document.
% es-noshorthands: disable babel-spanish's active `"`, `~`, `...` shorthands,
% which otherwise clash with biblatex-apa's \AddToHook{cmd/@makecaption/before}.
\usepackage[spanish, es-noshorthands, main=english]{babel}
\usepackage{csquotes}             % required by biblatex-apa

% --- Typography ------------------------------------------------------------
\usepackage{setspace}
\setstretch{1.15}                 % ~1.15 line spacing (justified is the default)
\usepackage{fancyhdr}             % page number top-right

% --- Figures ---------------------------------------------------------------
\usepackage{graphicx}
\graphicspath{{figures/}}         % look for PDF/PNG figures in figures/

% --- References: APA 7 via biblatex/biber ----------------------------------
\usepackage[style=apa, backend=biber]{biblatex}

% --- Hyperlinks (load last) -------------------------------------------------
\usepackage[hidelinks]{hyperref}

% ---------------------------------------------------------------------------
%  upb-strings.tex — Institutional strings of the UPB template.
%
%  SINGLE SOURCE OF TRUTH for every fixed institutional string.
%  Do NOT hardcode these strings anywhere else (especially not in chapters/).
%  To produce an English-only front matter, swap/override the Spanish macros
%  here — translation of the institutional boilerplate is a single-file change.
%
%  NOTE (faithful to the source PDF): the source writes "TITULO" without the
%  accent and keeps the faculty value in Spanish on the English page. Those
%  are reproduced here as-is where noted; adjust if the institution updates it.
% ---------------------------------------------------------------------------

% ---- Cover / portada -------------------------------------------------------
\newcommand{\fcTitleEs}{Diseño e Implementación de la Plataforma Fintcart para el Fortalecimiento de la Educación Financiera Digital}
\newcommand{\fcTitleEn}{Design and Implementation of the Fintcart Platform for Strengthening Digital Financial Education}
\newcommand{\fcAuthor}{Santiago Molano García}
\newcommand{\fcTrabajoGrado}{Trabajo de grado presentada para optar al título de}
\newcommand{\fcTituloOtorgado}{Ingeniero de Sistemas e Informática}
\newcommand{\fcDirectorLabel}{Director}
\newcommand{\fcDirectorName}{Mgtr. Danith Patricia Solórzano Escobar}
\newcommand{\fcUniversidad}{Universidad Pontificia Bolivariana}
\newcommand{\fcEscuela}{Escuela de Ingenierías}
\newcommand{\fcPrograma}{Ingeniería de Sistemas e Informática}
\newcommand{\fcCiudad}{Bucaramanga}
\newcommand{\fcAnio}{2026}

% ---- Front matter headings / lists -----------------------------------------
\newcommand{\fcContenido}{Contenido}
\newcommand{\fcListaTablas}{Lista de tablas}
\newcommand{\fcListaFiguras}{Lista de figuras}

% ---- Resumen (Spanish) -----------------------------------------------------
\newcommand{\fcResumenGeneralEs}{Resumen general de trabajo de grado en español}
\newcommand{\fcTituloLabel}{TITULO}           % source PDF writes it without accent
\newcommand{\fcAutorLabel}{AUTOR(ES)}
\newcommand{\fcProgramaLabel}{PROGRAMA}
\newcommand{\fcDirectorLabelEs}{DIRECTOR(A)}
\newcommand{\fcResumenLabel}{RESUMEN}
\newcommand{\fcPalabrasClaveLabel}{PALABRAS CLAVE}
\newcommand{\fcProgramaResumen}{Ing. Sistemas e Informática}

% TODO: replace these placeholder abstract/keywords with the final content.
\newcommand{\fcResumenEsTexto}{Este proyecto desarrolla una plataforma interactiva de educación financiera para abordar la carencia de conocimientos financieros en los métodos educativos modernos. La solución emplea una arquitectura de microservicios con tecnologías modernas para ofrecer contenido educativo escalable y atractivo mediante cuestionarios, simuladores y artículos informativos.}
\newcommand{\fcPalabrasClaveEsTexto}{Educación financiera, inclusión financiera, microservicios, plataforma web, simuladores.}
\newcommand{\fcVoBoEs}{V° B° DIRECTOR DE TRABAJO DE GRADO}

% ---- Summary (English) -----------------------------------------------------
\newcommand{\fcSummaryGeneralEn}{General summary of work of grade}
\newcommand{\fcTitleLabelEn}{TITLE}
\newcommand{\fcAuthorLabelEn}{AUTHOR(S)}
\newcommand{\fcFacultyLabelEn}{FACULTY}
\newcommand{\fcDirectorLabelEn}{DIRECTOR}
\newcommand{\fcAbstractLabelEn}{ABSTRACT}
\newcommand{\fcKeywordsLabelEn}{KEYWORDS}
% Source PDF leaves this value in Spanish on the English page; translated here.
\newcommand{\fcFacultyResumen}{Systems and Informatics Engineering}

% TODO: replace these placeholder abstract/keywords with the final content.
\newcommand{\fcAbstractEnTexto}{This project develops an interactive financial education platform to address the gap in financial literacy in modern educational methods. The solution employs a microservices architecture with modern technologies to deliver scalable and engaging educational content through quizzes, simulators, and informative articles.}
\newcommand{\fcKeywordsEnTexto}{Financial education, financial inclusion, microservices, web platform, simulators.}
\newcommand{\fcVoBoEn}{V° B° DIRECTOR OF GRADUATE WORK}

% ---------------------------------------------------------------------------
%  commands.tex — Shared macros (layout logic only).
%  Institutional *strings* live in preamble/upb-strings.tex — never here.
% ---------------------------------------------------------------------------

% --- Optional lists of tables/figures ---------------------------------------
% Flip to \printliststrue once the document has at least one table/figure.
% Empty lists are never printed (and never error).
\newif\ifprintlists
\printlistsfalse

% --- Monospace for figures / money / code -----------------------------------
% N-15: a monetary figure must never be truncated — mono keeps digits stable.
\newcommand{\fcnum}[1]{\texttt{#1}}   % money / numeric literals (may switch to siunitx later)
\newcommand{\code}[1]{\texttt{#1}}    % inline code / identifiers

% --- Page style: page number top-right on every page -------------------------
\pagestyle{fancy}
\fancyhf{}
\fancyhead[R]{\thepage}
\renewcommand{\headrulewidth}{0pt}
\renewcommand{\footrulewidth}{0pt}

% --- Depth of the table of contents ------------------------------------------
\setcounter{tocdepth}{2}

% --- Portada (cover page) -----------------------------------------------------
% #1 = optional director block (empty on cover page 1).
\newcommand{\fcportada}[1]{%
  \null\vfill
  \begin{center}
    {\Large\bfseries\fcTitleEs}\par
    \vspace{1.6cm}
    {\Large\fcAuthor}\par
    \vspace{1.6cm}
    {\normalsize\fcTrabajoGrado}\par
    {\normalsize\bfseries\fcTituloOtorgado}\par
    \vspace{1.6cm}
    #1
    \vspace{1.6cm}
    {\normalsize\fcUniversidad}\par
    {\normalsize\fcEscuela}\par
    {\normalsize\fcPrograma}\par
    {\normalsize\fcCiudad}\par
    {\normalsize\fcAnio}\par
  \end{center}
  \vfill\null
  \clearpage
}

% Director block for cover page 2.
\newcommand{\fcDirectorBlock}{%
  {\normalsize\fcDirectorLabel}\par
  {\normalsize\fcDirectorName}\par
}

% --- Resumen general de trabajo de grado (Spanish) ----------------------------
\newcommand{\fcresumenes}{%
  \begin{center}
    {\large\bfseries\fcResumenGeneralEs}\par
  \end{center}
  \vspace{1em}
  \noindent
  \begin{tabular}{@{}l@{\hspace{2em}}p{11.5cm}@{}}
    \fcTituloLabel:     & \fcTitleEs \\
    \fcAutorLabel:      & \fcAuthor \\
    \fcProgramaLabel:   & \fcProgramaResumen \\
    \fcDirectorLabelEs: & \fcDirectorName \\
  \end{tabular}

  \vspace{2em}
  \begin{center}{\large\bfseries\fcResumenLabel}\end{center}
  \noindent \fcResumenEsTexto

  \vspace{1.5em}
  \noindent \fcPalabrasClaveLabel:\par
  \vspace{0.5em}
  \noindent \fcPalabrasClaveEsTexto

  \vspace{2em}
  \noindent \fcVoBoEs
  \clearpage
}

% --- General summary of work of grade (English) -------------------------------
\newcommand{\fcsummaryen}{%
  \begin{center}
    {\large\bfseries\fcSummaryGeneralEn}\par
  \end{center}
  \vspace{1em}
  \noindent
  \begin{tabular}{@{}l@{\hspace{2em}}p{11.5cm}@{}}
    \fcTitleLabelEn:    & \fcTitleEn \\
    \fcAuthorLabelEn:   & \fcAuthor \\
    \fcFacultyLabelEn:  & \fcFacultyResumen \\
    \fcDirectorLabelEn: & \fcDirectorName \\
  \end{tabular}

  \vspace{2em}
  \begin{center}{\large\bfseries\fcAbstractLabelEn}\end{center}
  \noindent \fcAbstractEnTexto

  \vspace{1.5em}
  \noindent \fcKeywordsLabelEn:\par
  \vspace{0.5em}
  \noindent \fcKeywordsEnTexto

  \vspace{2em}
  \noindent \fcVoBoEn
  \clearpage
}


\addbibresource{refs/referencias.bib}

% Compile only a subset of chapters (fast iteration). Uncomment and list the
% files you want; all others are skipped but keep their section numbers:
% \includeonly{chapters/03-objetivos}

\begin{document}

% ============================ FRONT MATTER (Spanish) ============================
\selectlanguage{spanish}
\renewcommand{\contentsname}{\fcContenido}
\renewcommand{\listtablename}{\fcListaTablas}
\renewcommand{\listfigurename}{\fcListaFiguras}

% --- 1. Portada (cover page 1) ---
\fcportada{}

% --- 2. Portada (cover page 2, with director) ---
\fcportada{\fcDirectorBlock}

% --- 3. Contenido (table of contents) ---
\tableofcontents
\clearpage

% --- 4. Resumen general de trabajo de grado en español ---
\fcresumenes

% --- 5. General summary of work of grade (English) ---
\fcsummaryen

% --- 6. Lists of tables and figures (only if non-empty) ---
\ifprintlists
  \listoftables
  \clearpage
  \listoffigures
  \clearpage
\fi

% ============================ BODY (English) ============================
\selectlanguage{english}

% TRANSLATION UNIT: translate this file only.
% ===========================================================================
%  01 — Planteamiento del problema
%  The \section title is the canonical (Spanish) chapter name; the body prose
%  is written in ENGLISH and translated to Spanish later.
% ===========================================================================

\section{Planteamiento del problema}

% TODO: Problem statement (English source text).
Describe the problem this platform addresses: the gap in financial literacy in
modern educational methods and the lack of accessible, interactive tools for
financial education in the Colombian context.

% TODO: Reserve a causes/effects tree figure here (required by the rubric):
%   \begin{figure}[htbp]
%     \centering
%     \includegraphics[width=\linewidth]{figures/arbol-causas-efectos.pdf}
%     \caption{Causes and effects tree of the problem.}
%     \label{fig:arbol-causas-efectos}
%   \end{figure}

% TRANSLATION UNIT: translate this file only.
% ===========================================================================
%  02 — Justificación
%  The \section title is the canonical (Spanish) chapter name; the body prose
%  is written in ENGLISH and translated to Spanish later.
% ===========================================================================

\section{Justificación}

% TODO: Justification (English source text).
Explain why this project matters: the importance of financial education for
social and individual economic well-being, and why a digital platform is an
appropriate vehicle for it.

% TRANSLATION UNIT: translate this file only.
% ===========================================================================
%  03 — Objetivos
%  The \section title is the canonical (Spanish) chapter name; the body prose
%  is written in ENGLISH and translated to Spanish later.
%  Labels OE1..OE5 are kept ready so other chapters can \ref them.
% ===========================================================================

\section{Objetivos}

\subsection{Objetivo general}

% TODO: General objective (English source text).
State the single general objective of the project.

\subsection{Objetivos específicos}

\begin{itemize}
  \item[OE1.] \label{oe1} % TODO: specific objective 1 (English).
  \item[OE2.] \label{oe2} % TODO: specific objective 2 (English).
  \item[OE3.] \label{oe3} % TODO: specific objective 3 (English).
  \item[OE4.] \label{oe4} % TODO: specific objective 4 (English).
  \item[OE5.] \label{oe5} % TODO: specific objective 5 (English).
\end{itemize}

% TRANSLATION UNIT: translate this file only.
% ===========================================================================
%  04 — Marco referencial
%  The \section title is the canonical (Spanish) chapter name; the body prose
%  is written in ENGLISH and translated to Spanish later.
% ===========================================================================

\section{Marco referencial}

% TODO: Referential/theoretical framework (English source text), citing from
% refs/referencias.bib with \parencite{} / \textcite{}.
Cover the theoretical background: financial education, financial mathematics
regimes (nominal vs. effective rates, French amortization, simple vs. compound
interest) and their Colombian regulatory anchors (UVT, GMF, usury rate, IPC).

% TRANSLATION UNIT: translate this file only.
% ===========================================================================
%  05 — Metodología
%  The \section title is the canonical (Spanish) chapter name; the body prose
%  is written in ENGLISH and translated to Spanish later.
% ===========================================================================

\section{Metodología}

% TODO: Methodology (English source text).
Declare the hybrid approach: Kanban (flow + WIP limits) + XP (CI, automated
testing, continuous refactoring) inside a document-driven structure, executed
as 4 increments with verifiable acceptance criteria. Include phases and the
objectives-to-development traceability matrix.

% TRANSLATION UNIT: translate this file only.
% ===========================================================================
%  06 — Requerimientos
%  The \section title is the canonical (Spanish) chapter name; the body prose
%  is written in ENGLISH and translated to Spanish later.
% ===========================================================================

\section{Requerimientos}

% TODO: Requirements (English source text).
Structure the functional and non-functional requirements (FR/NFR) of the three
specs plus the six docs/rf/RF-*.md documents, in an IEEE-830 style.

% TRANSLATION UNIT: translate this file only.
% ===========================================================================
%  07 — Arquitectura del sistema
%  The \section title is the canonical (Spanish) chapter name; the body prose
%  is written in ENGLISH and translated to Spanish later.
% ===========================================================================

\section{Arquitectura del sistema}

% TODO: System architecture (English source text).
Document the polyglot microservices architecture: C4 L1/L2/L3 diagrams,
containers and components, gRPC + REST-on-the-edge, PostgreSQL database-per-
service, Redis and RabbitMQ boundaries, plus a deployment diagram. Justify
technology decisions (promote D-13..D-32 to ADRs).

% TRANSLATION UNIT: translate this file only.
% ===========================================================================
%  08 — Modelo de datos
%  The \section title is the canonical (Spanish) chapter name; the body prose
%  is written in ENGLISH and translated to Spanish later.
% ===========================================================================

\section{Modelo de datos}

% TODO: Data model (English source text).
Document the data model per service (data-model.md) and one ER diagram per
database, covering the persistent schema of the eight services.

% TRANSLATION UNIT: translate this file only.
% ===========================================================================
%  09 — Diseño de interfaz y experiencia
%  The \section title is the canonical (Spanish) chapter name; the body prose
%  is written in ENGLISH and translated to Spanish later.
% ===========================================================================

\section{Diseño de interfaz y experiencia}

% TODO: UI/UX design (English source text).
Document the design system: tokens, 19 primitives, accessibility (AA contrast,
minimum 360 px), the four breakpoint tokens (480/768/1024/1280) and the
benchmark of reference products.

% TRANSLATION UNIT: translate this file only.
% ===========================================================================
%  10 — Implementación
%  The \section title is the canonical (Spanish) chapter name; the body prose
%  is written in ENGLISH and translated to Spanish later.
% ===========================================================================

\section{Implementación}

% TODO: Implementation (English source text).
Document the implementation of the eight services and the formula engine
(parser \textrightarrow\ persisted AST in JSONB \textrightarrow\ bounded evaluator in
rust\_decimal), including the seed calculator definitions.

% TRANSLATION UNIT: translate this file only.
% ===========================================================================
%  11 — Pruebas
%  The \section title is the canonical (Spanish) chapter name; the body prose
%  is written in ENGLISH and translated to Spanish later.
% ===========================================================================

\section{Pruebas}

% TODO: Tests (English source text).
Document the test suites: unit, integration, end-to-end (Playwright) and the
seed\_regression suite that guards the formula engine before deleting the
legacy native calculators.

% TRANSLATION UNIT: translate this file only.
% ===========================================================================
%  12 — Despliegue
%  The \section title is the canonical (Spanish) chapter name; the body prose
%  is written in ENGLISH and translated to Spanish later.
% ===========================================================================

\section{Despliegue}

% TODO: Deployment (English source text).
Document the real deployment: VPS + Caddy + TLS, Docker Compose, Kubernetes
manifests and CI/CD.

% TRANSLATION UNIT: translate this file only.
% ===========================================================================
%  13 — Evaluación del sistema
%  The \section title is the canonical (Spanish) chapter name; the body prose
%  is written in ENGLISH and translated to Spanish later.
% ===========================================================================

\section{Evaluación del sistema}

% TODO: System evaluation (English source text).
Document the evaluation evidence: Lighthouse across breakpoints, k6 load
testing, usability sessions and the System Usability Scale (SUS).

% TRANSLATION UNIT: translate this file only.
% ===========================================================================
%  14 — Conclusiones y trabajo futuro
%  The \section title is the canonical (Spanish) chapter name; the body prose
%  is written in ENGLISH and translated to Spanish later.
% ===========================================================================

\section{Conclusiones y trabajo futuro}

% TODO: Conclusions and future work (English source text).
Summarize the results against the objectives (OE1..OE5) and list future work.

% TRANSLATION UNIT: translate this file only.
% ===========================================================================
%  15 — Referencias
%  The \section title is the canonical (Spanish) chapter name.
% ===========================================================================

\section{Referencias}

% DRAFTING: while chapters are still stubs nothing is cited yet, so biblatex would
% print an empty list. \nocite{*} renders every entry in refs/referencias.bib so the
% bibliography can be proofread. It is harmless once real citations exist (cited
% works are printed anyway) — but remove it before delivery if a reviewer objects
% to uncited entries appearing in the list.
\nocite{*}

% Prints the bibliography from refs/referencias.bib (APA 7 via biblatex).
% heading=none avoids a second "References" heading under the section title.
\printbibliography[heading=none]


% ============================ APPENDICES ============================
% \appendix is issued inside 16-anexos.tex; everything after is lettered A, B, …

% TRANSLATION UNIT: translate this file only.
% ===========================================================================
%  16 — Anexos (appendix section)
%  \appendix switches subsequent \section numbering to letters (A, B, ...).
%  The \section title is the canonical (Spanish) chapter name; the body prose
%  is written in ENGLISH and translated to Spanish later.
% ===========================================================================

\appendix
\section{Anexos}
\label{anex:anexos}

% TODO: Appendices (English source text). Add each annex as a subsection or a
%   lettered \section, e.g.:
%   \subsection{Deployment manifests}
%   \subsection{API contracts (OpenAPI + protos)}
%   \subsection{Schedule / timeline}


\end{document}
